This chapter gives an introduction to the project report.

5-6 pages with a full summary


The NZ population is generally one of the most skeptical and concerned about AI in the developed world \cite{TrustAttitudesUse,GlobalAttitudesAI2026}. We are also a country with some of the weakest AI governance and oversight \cn{find reference about feather tough approach, we weren't at various AI safety conferences, have no AISI...}


\section{Report overview}

Following on from this introduction chapter there are 5 more chapters to the report. The next chapter \ref{chap:background} gives the key backgroun

This report is structured into 3 main chapters, each focusing on a specific part of the project. For each chapter of the project it explains the methods used and the results obtained. Each chapter is self contained and can be read somewhat independently of the others. The 3 main chapters are:

\begin{itemize}
    \item \ref{chap:finetuning}: Fine-tuning of the LLM on WVS data
    \item \ref{chap:simulation}: Behavioural analysis of the fine-tuned and base LLMs
    \item \ref{chap:survey}: Public consultation on the alignment of the LLMs
\end{itemize}

The final chapter \ref{chap:conclusion} concludes the report with a summary of the findings, conclusions drawn from the project, key limitations and recommendations for future work.

\section{Research questions}

This project seeks to answer two research questions:

\begin{itemize}
    \item \newrq{align}{Can we use publicly collected value survey data to modify the behaviour of LLMs so that they align with the values of the public?}
    \item \newrq{prefer}{Will users prefer the behaviour of models aligned to their values over unaligned base models?}
\end{itemize}
